\chapter*{\large ЗАКЛЮЧЕНИЕ}
\addcontentsline{toc}{chapter}{ЗАКЛЮЧЕНИЕ}

В ходе выполнения курсовой работы были изучены основные алгоритмы процедурной генерации изображений и псевдослучайных структур: шум Перлина, симплексный шум, алгоритм Diamond-Square и клеточные автоматы. Были рассмотрены принципы работы каждого метода, их математическая основа и особенности применения при генерации карт высот, текстур и бинарных структур.

В практической части работы были реализованы алгоритмы генерации на языке Python. Для визуализации результатов было разработано приложение с графическим интерфейсом на базе библиотеки Tkinter. Отображение матриц выполняется с помощью Matplotlib, что позволяет представлять результат работы алгоритмов как полутоновое изображение. В приложении предусмотрены окна настройки параметров: размера области, масштаба, seed, количества октав, параметров клеточного автомата, а также параметров алгоритма Diamond-Square.

В результате работы были реализованы следующие возможности:
\begin{enumerate}
    \item генерация шума Перлина и симплексного шума;
    \item генерация фрактальной карты высот алгоритмом Diamond-Square;
    \item применение клеточного автомата к бинарной матрице;
    \item построение многооктавного шума за счёт сложения слоёв с увеличивающейся частотой и уменьшающейся амплитудой;
    \item комбинирование нескольких шумовых карт;
    \item применение маски острова и бинаризация результата по заданному порогу;
    \item отображение и сохранение результата в виде изображения.
\end{enumerate}

Шум Перлина показал себя как удобный способ получения плавных переходов и естественных текстур. Симплексный шум оказался более универсальным при работе с пространствами большей размерности. Алгоритм Diamond-Square удобен для генерации карт высот с фрактальной структурой. Клеточные автоматы, в свою очередь, позволяют преобразовывать бинарные карты и получать структуры, похожие на пещеры, области суши и воды.

Таким образом, цель курсовой работы была достигнута: были изучены алгоритмы процедурной генерации, выполнена их программная реализация и создано приложение с графическим интерфейсом для визуализации и комбинирования результатов. Полученная программа может быть расширена в дальнейшем: например, можно добавить цветовые карты биомов, экспорт высотных карт, трёхмерную визуализацию рельефа, ускорение вычислений с помощью NumPy или Numba, а также более гибкую систему слоёв и масок.
